% TOP_PROBLEM: {"id": 1153, "title": "Polignac's Conjecture", "category_id": 1, "category": {"id": 1, "name": "number_theory", "display_name": "Number Theory", "description": "Properties of integers, prime numbers, Diophantine equations.", "slug": "number-theory", "order_index": 1, "created_at": "2026-07-31T15:26:25.670Z"}, "status": "open"}
% FIRST_POSTED: 2026-09-27
% !TeX program = lualatex
% Compile twice: lualatex 1153.tex
% All references and prose are embedded; no BibTeX database or external inputs.
\documentclass[11pt,a4paper]{article}
\usepackage[margin=24mm,headheight=14pt]{geometry}
\usepackage{fontspec}
\setmainfont{Latin Modern Roman}
\setsansfont{Latin Modern Sans}
\setmonofont{Latin Modern Mono}
\usepackage{amsmath,amssymb,mathtools}
\usepackage{unicode-math}
\setmathfont{Latin Modern Math}
\usepackage{microtype}
\usepackage{enumitem,xurl,needspace}
\usepackage[dvipsnames]{xcolor}
\usepackage{hyperref}
\hypersetup{unicode=true,colorlinks=true,linkcolor=MidnightBlue,urlcolor=MidnightBlue,
 pdftitle={Research attempt: Problem 1153},pdfauthor={Research notebook},bookmarksnumbered=true}
\usepackage{bookmark}
\usepackage{tocloft}
\setlength{\cftsecnumwidth}{3em}
\cftsetpnumwidth{2.5em}
\cftsetrmarg{4em}
\usepackage{titlesec}
\titleformat{\section}{\Large\bfseries\raggedright}{\thesection}{0.7em}{}
\usepackage{fancyhdr}
\pagestyle{fancy}\fancyhf{}
\fancyhead[L]{\small\sffamily Open Problems Math | Research notebook}
\fancyhead[R]{\small Problem 1153 | 27 September 2026}
\fancyfoot[C]{\thepage}
\setlength{\parindent}{0pt}
\setlength{\parskip}{3pt plus 1pt}
\setlength{\emergencystretch}{3em}
\setcounter{tocdepth}{1}
\setcounter{secnumdepth}{1}
\setlist[itemize]{leftmargin=3em,itemsep=5pt,topsep=4pt,parsep=0pt}
\allowdisplaybreaks
\newcommand{\N}{\mathbb{N}}
\newcommand{\Z}{\mathbb{Z}}
\newcommand{\Q}{\mathbb{Q}}
\newcommand{\R}{\mathbb{R}}
\newcommand{\C}{\mathbb{C}}
\newcommand{\F}{\mathbb{F}}
\newcommand{\A}{\mathbb{A}}
\newcommand{\PP}{\mathbb P}
\newcommand{\E}{\mathbb{E}}
\newcommand{\eps}{\varepsilon}
\newcommand{\dd}{\,\mathrm d}
\newcommand{\sref}[2]{\hyperlink{source:#1:#2}{[S#2]}}
\newcommand{\eref}[2]{\hyperlink{extra:#1:#2}{[E#2]}}
% Turn the next switch off for a shorter reading copy. The quotations preserve
% every exactTarget field, including qualifications omitted from a short summary.
\newif\ifshowcatalogue
\showcataloguetrue
\newcommand{\cataloguescope}[1]{\ifshowcatalogue
 \paragraph{Exact scope in the supplied catalogue.}
 \begin{quote}\small #1\end{quote}\fi}
\usepackage{etoolbox}
\AtBeginEnvironment{itemize}{\small}
\numberwithin{equation}{section}
% Standalone export. Original research content preserved.
\begin{document}
\setcounter{section}{0}
% ================================================================
% problemId: 1153
% Canonical problem ID: 1153
% exactTarget SHA-256: 0b6d706e31b161de6b8ada3302e573f70856be4acb967c0dd62af7e65cdcf175
\clearpage
\section*{\texorpdfstring{Polignac's Conjecture}{Polignac's Conjecture}}\label{1153}
\subsection{Definitions and mathematical statement}
Write $p_1<p_2<\cdots$ for the increasing sequence of primes. Polignac's assertion is
\[
 \forall k\in\Z_{>0}\ \forall N\ \exists n>N:\qquad p_{n+1}-p_n=2k.
\]
Consecutiveness means there is no other prime between the pair. Finding infinitely many pairs of primes at distance $2k$ without this condition would be a different and potentially weaker assertion for $k>1$.
\par\smallskip\textit{Formulation and concept references:} \sref{1153}{1}.
\cataloguescope{For every positive even integer 2k, there are infinitely many pairs of consecutive primes whose difference is 2k.}
\subsection{Short English statement}
Choose any positive even gap. Must that exact spacing recur between neighboring primes arbitrarily far along the prime sequence?
% BEGIN RESEARCH ADDITION 106
\subsection{Research attempt}

\paragraph{Current frontier.}
The established bounded-gap theorem gives infinitely many neighboring prime
pairs with a gap at most $246$; it does not specify each even gap. A preprint
of Stadlmann, posted on 31 August 2026, reports the improved bound $240$ by
combining distribution estimates. This recent result is recorded as a
preprint, not independently verified here, and is not represented as a
resolution of any prescribed-gap assertion. \eref{1153}{1}\eref{1153}{2}
Maynard's sieve explains how to find several primes among sufficiently many
admissible translates; it need not make two \emph{specified} translates
simultaneously prime. \eref{1153}{3}

\paragraph{Attack route.}
I separate three issues often conflated in this problem: obtaining prime
endpoints, eliminating interior primes, and making the endpoint argument work
for every fixed even gap. There are two ways to handle the middle issue.
One can show that interior-prime exceptions are rarer than the expected
endpoint pairs, or prearrange composite interior values using congruences.
I develop both reductions, then test whether the elementary two-form
Maynard variational problem can supply their missing endpoint input.
Neither reduction assumes consecutiveness at the outset.

\paragraph{Attempt.}
\textbf{1. An endpoint-to-consecutive-gap estimate.}
Fix an even $h\geq2$. For $X>h$, define
\begin{align*}
 P_h(X)&=\#\{p:\ p,p+h\text{ prime},\ p+h\leq X\},\\
 G_h(X)&=\#\{p:\ p,p+h\text{ consecutive primes},\ p+h\leq X\},\\
 T_{h,a}(X)&=\#\{p:\ p,p+a,p+h\text{ prime},\ p+h\leq X\}
 \quad(1\leq a<h).
\end{align*}
An endpoint pair is excluded from $G_h$ only if some intermediate translate
is prime. The union bound therefore gives the exact inequalities
\begin{equation}\label{1153:union}
 0\leq P_h(X)-G_h(X)\leq\sum_{a=1}^{h-1}T_{h,a}(X).
\end{equation}
For each fixed $a$, a three-dimensional upper-bound sieve gives
\begin{equation}\label{1153:triple}
 T_{h,a}(X)\ll_h \frac{X}{(\log X)^3}.
\end{equation}
Here are the hypothesis checks. If $\{0,a,h\}$ is admissible, the three
linear forms have nonzero discriminant factor $ah(h-a)$ and occupy fewer
than $\ell$ residue classes for every prime $\ell$. Apply the standard
prime-tuple upper bound, for example the $K=3$ instance of Theorem 1 in
\eref{1153}{4}. If the triple is inadmissible, choose a prime $\ell$ at which
its residues cover every class. One of $p,p+a,p+h$ must then equal $\ell$,
so there are at most three candidate values of $p$. Thus this case is bounded
as well. The number of shifts is fixed with $h$; no uniformity in growing
$h$ is being claimed.

Consequently there is a constant $C_h$ such that for all sufficiently large
$X$,
\begin{equation}\label{1153:gapcount}
 P_h(X)-C_h\frac{X}{(\log X)^3}\leq G_h(X)\leq P_h(X).
\end{equation}
This estimate isolates a quantitative input sufficient for the selected
problem:
\begin{equation}\label{1153:pairinput}
 \text{for every fixed even $h\geq2$,}\qquad
 \limsup_{X\to\infty}\frac{P_h(X)(\log X)^3}{X}=\infty.
\end{equation}
Indeed choose arbitrarily large $X$ for which the ratio exceeds $C_h+1$.
Then $G_h(X)\geq X/(\log X)^3\to\infty$, which gives infinitely many
consecutive gaps $h$. Any lower bound $P_h(X)\gg_h X/(\log X)^2$ along an
unbounded sequence of $X$ would suffice. Mere infinitude of $P_h$ would
\emph{not}: an infinite but sufficiently sparse set could fit entirely
inside the exception bound in \eqref{1153:gapcount}.

For the edge cases, $h=2$ permits no intermediate prime, so $P_2=G_2$.
For $h=4$, the only nonconsecutive prime pair is $(3,7)$: otherwise an
interior prime would produce three primes $p,p+2,p+4$, one divisible by $3$.
Thus $P_4(X)-G_4(X)=1$ for $X\geq7$. The general quantitative input is a
sufficient statement, not a claim that these easy edge cases require it.

\medskip
\textbf{2. Congruence preconditioning with no density assumption.}
A different reduction eliminates the interior primes in advance. For every
even integer $a$ with $2\leq a\leq h-2$, choose a distinct prime $q_a>h$.
Set
\[
 M=2\prod_{a\in\{2,4,\ldots,h-2\}}q_a.
\]
By the Chinese remainder theorem choose $0\leq r<M$ with
\begin{equation}\label{1153:crt}
 r\equiv1\pmod2,\qquad r\equiv-a\pmod{q_a}.
\end{equation}
The two linear forms
\[
 L_0(t)=Mt+r,\qquad L_h(t)=Mt+r+h
\]
are jointly admissible. If $\ell\mid M$, neither endpoint vanishes modulo
$\ell$: for $\ell=q_a$ their residues are $-a$ and $h-a$, both nonzero
because $0<a<h<q_a$; at $2$ both are odd. If $\ell\nmid M$, the two forms
exclude at most two classes modulo $\ell$, and $\ell\geq3$. This also proves
the primitivity of both forms.

If both endpoints are prime and $t$ is large enough, each interior value
is composite. Odd $a$ gives an even integer larger than $2$. Even $a$ gives
a multiple of its assigned $q_a$ larger than $q_a$. Therefore:
\begin{quote}
Infinitely many simultaneous prime values of $L_0,L_h$ imply infinitely
many \emph{consecutive} gaps of size $h$.
\end{quote}
This is an unconditional implication with an unproved simultaneous-primality
hypothesis. In the language of prime tuples, the two-form Dickson assertion
for these explicit progressions would suffice. No asymptotic count is needed
for this second reduction. The prime-tuple framework and its distinction
from the presently known many-translates sieve are discussed in
\eref{1153}{3}\eref{1153}{4}.

For a concrete check, at $h=6$ choose $q_2=7$, $q_4=11$. Then $M=154$ and
$r=117$ work. At $t=1$ the endpoints are $271$ and $277$; their five interior
values are composite, with $273$ divisible by $7$ and $275$ by $11$.
One such instance checks the construction, not its infinitude.

There is also no hidden vanishing singular series. Writing $\nu_\ell$ for
the number of roots of $L_0L_h$ modulo $\ell$, each local factor
\[
 (1-\nu_\ell/\ell)(1-1/\ell)^{-2}
\]
is positive. Outside finitely many primes $\nu_\ell=2$, and the factor is
$1+O(\ell^{-2})$; hence their infinite product converges to a positive
number. Positivity is a consistency check on a proposed asymptotic, not a
proof that either endpoint is prime infinitely often together with the other.

\medskip
\textbf{3. A sharper obstruction in the basic two-form sieve.}
Let $\mathcal R=\{(s,t):s,t\geq0,\ s+t\leq1\}$ and let $F\in L^2(\mathcal R)$
be nonzero. Extend it by zero off $\mathcal R$, and write
\begin{align*}
 I(F)&=\iint_{\mathcal R}|F(s,t)|^2\,ds\,dt,\\
 J_1(F)&=\int_0^1\left|\int_0^{1-t}F(s,t)\,ds\right|^2dt,\\
 J_2(F)&=\int_0^1\left|\int_0^{1-s}F(s,t)\,dt\right|^2ds.
\end{align*}
I claim the uniform estimate
\begin{equation}\label{1153:M2}
 J_1(F)+J_2(F)\leq2\log2\,I(F).
\end{equation}
This is more informative than bounding each $J_i$ separately by $I$.

\emph{Proof.} Put
\[
 w_1(s,t)=\frac{1+s-t}{2},\qquad
 w_2(s,t)=\frac{1+t-s}{2}.
\]
They are positive in the interior and satisfy $w_1+w_2=1$. For $t<1$,
\[
 \int_0^{1-t}\frac{ds}{w_1(s,t)}=2\log2.
\]
Weighted Cauchy--Schwarz therefore gives
\[
 \left|\int_0^{1-t}F(s,t)\,ds\right|^2
 \leq2\log2\int_0^{1-t}w_1(s,t)|F(s,t)|^2\,ds.
\]
Integrate in $t$ and repeat with $s,t$ interchanged. Their sum is
\eqref{1153:M2}; the boundary where a weight vanishes has measure zero.
All integrands on the right are integrable because $0\leq w_i\leq1$.
$\square$

In Maynard's basic simplex-supported setup, a level of distribution $\theta$
yields the leading weighted prime-count ratio
\[
 \frac{\theta}{2}\frac{J_1(F)+J_2(F)}{I(F)}.
\]
For two forms, \eqref{1153:M2} bounds it by $\theta\log2<1$ even if
$\theta=1$. To force both forms prime by the nonnegative weighted-count
criterion, that ratio must exceed $1$. Thus this particular two-form mean
criterion cannot supply the missing input, even with a full level of
distribution. The factor $\theta/2$, the support restriction, and the
threshold are exactly those of Proposition 4.2 and its proof in
\eref{1153}{3}. This does not exclude different weights, different support
geometries, additional correlations, or genuinely different sieve arguments.
No novelty claim is made for the variational inequality.

For example $F=1$ on $\mathcal R$ gives $I=1/2$, $J_1=J_2=1/3$, hence
$(J_1+J_2)/I=4/3<2\log2$. Increasing the quality of this test function within
the stated class cannot cross the required threshold.

\paragraph{Stress test.}
The following exact enumeration, reproduced by the companion script, counts
all endpoint pairs and neighboring gaps with upper endpoint at most
$10^6$. The last column counts interior-prime incidences, with multiplicity,
so it is the right side of \eqref{1153:union}, not the number of bad pairs.
\begin{center}
\begin{tabular}{r|r|r|r}
$h$ & $P_h(10^6)$ & $G_h(10^6)$ & $\sum_a T_{h,a}(10^6)$\\\hline
2 & 8169 & 8169 & 0\\
4 & 8144 & 8143 & 1\\
6 & 16386 & 13549 & 2837\\
10 & 10934 & 7079 & 4174\\
12 & 16378 & 8005 & 10006
\end{tabular}
\end{center}
These finite counts test the definitions and the union-bound direction;
they say nothing about the limit in \eqref{1153:pairinput}. In particular,
the endpoint-pair and consecutive-gap columns are not interchangeable.

The CRT modulus grows with the chosen $h$, but that is allowed: Polignac
quantifies over each \emph{fixed} gap and does not demand estimates uniform
in $h$. The sufficiently-large-$t$ qualification excludes a forced divisor
being equal to the interior value itself. A theorem producing infinitely
many prime values for each individual linear form is not enough, since the
intersection of two infinite sets can be empty. Likewise, Bunyakovsky's
single-polynomial assertion from UnsolvedMath problem 1292 cannot simply be applied to
$L_0L_h$: that product is reducible. Both reductions still require a genuine
two-prime correlation result, and at $h=2$ they necessarily include the twin
prime problem.

\paragraph{Outcome.}
\textbf{Outcome: Conditional reduction.}

Two complete implications have been established: the quantitative pair input
\eqref{1153:pairinput} implies all the prescribed consecutive gaps, and an
infinitude assertion for two explicit admissible arithmetic progressions
implies any one prescribed gap without requiring a density. The elementary
bound \eqref{1153:M2} explains why the most direct two-form version of the
basic Maynard criterion cannot provide that input. Neither endpoint
hypothesis has been proved. Thus this attempt has isolated consecutiveness
rigorously, but has not resolved Polignac's conjecture.
% END RESEARCH ADDITION 106
\subsection{Sources}
\begin{itemize}\raggedright
\item[\textnormal{[S1]}]\hypertarget{source:1153:1}{} MathOverflow contributors, What is the importance of Polignac's Conjecture? (2022).
\url{https://mathoverflow.net/questions/416389/what-is-the-importance-of-polignac-s-conjecture}.
{\small\textit{Catalogue formulation source.}}
% BEGIN REFERENCE ADDITION 106
\item[\textnormal{[E1]}]\hypertarget{extra:1153:1}{} D. H. J. Polymath,
\emph{Variants of the Selberg sieve, and bounded intervals containing many
primes}, Research in the Mathematical Sciences 1 (2014), article 12;
arXiv:1407.4897, Introduction, the unconditional $H_1\leq246$ bound.
\url{https://arxiv.org/abs/1407.4897}.
{\small\textit{Consulted 27 September 2026.}}
\item[\textnormal{[E2]}]\hypertarget{extra:1153:2}{} J. Stadlmann,
\emph{Bounded gaps between primes}, arXiv:2608.31126v1
(31 August 2026), Abstract, Introduction and the main $H_1\leq240$ theorem.
\url{https://arxiv.org/html/2608.31126v1}.
{\small\textit{Consulted 27 September 2026; recent preprint, not independently audited here.}}
\item[\textnormal{[E3]}]\hypertarget{extra:1153:3}{} J. Maynard,
\emph{Small gaps between primes}, Annals of Mathematics 181 (2015), 383--413;
arXiv:1311.4600, Section 4, Propositions 4.1--4.3 and equation (4.4).
\url{https://arxiv.org/html/1311.4600}.
{\small\textit{Consulted 27 September 2026; simplex support and the $\theta/2$ factor checked.}}
\item[\textnormal{[E4]}]\hypertarget{extra:1153:4}{} T. Dubbe,
\emph{Explicit bounds for prime $k$-tuplets}, arXiv:2409.04261v1
(6 September 2024), Theorem 1, especially its $K=3$ case and its
admissibility and nonzero-discriminant hypotheses.
\url{https://arxiv.org/html/2409.04261v1}.
{\small\textit{Consulted 27 September 2026; used for the standard upper-bound sieve estimate, not a prime-tuple lower bound.}}
% END REFERENCE ADDITION 106
\end{itemize}

\end{document}
